\section*{Chapter \ref{ch:m1:european-equity-market-structure} --- European Equity Market Structure}

\begin{solution}{exo:m1:european-equity-market-structure:1}
(a) \omterm{def:m1:european-equity-market-structure:rm}{Regulated market}. (b) \omterm{def:m1:exchanges-brokers-venues:mtf}{Multilateral trading facility}. (c) \omterm{def:m1:european-equity-market-structure:si}{Systematic
internaliser}. (d) \omterm{def:m1:european-equity-market-structure:rm}{Organised trading facility}, permitted because the
instruments are bonds.
\end{solution}

\begin{solution}{exo:m1:european-equity-market-structure:2}
$H = 0.2025 + 0.0625 + 0.0225 + 0.01 + 0.0025 = 0.30$; effective number
3.33.
\end{solution}

\begin{solution}{exo:m1:european-equity-market-structure:3}
Closing auctions $0.243 \times 57.5 = \text{\euro}14.0$ billion a day;
\omterm{def:m1:european-equity-market-structure:periodic}{periodic auctions} \euro 6.0 billion; non-displayed books \euro 6.3 billion.
Off-exchange $80 - 57.5 = 22.5$, of which \euro 14.3 billion \omterm{def:m1:european-equity-market-structure:si}{systematic
internalisers} and \euro 8.2 billion other.
\end{solution}

\begin{solution}{exo:m1:european-equity-market-structure:4}
Best bid 18.350 (Z), best offer 18.355 (Y). Buying at 18.355 on Y is at the
best offer: no issue this time. Had Z offered lower, the broker would still
have broken no \emph{market} rule --- there is no \omterm{def:m1:us-equity-market-structure:opr}{trade-through} prohibition
--- but it would have to justify under its best-execution policy why Z is
not among its venues (cost of connection against expected \omterm{def:m1:retail-flow-and-wholesaling:pfof}{price
improvement}). The obligation is about the policy and its review, not about
each fill.
\end{solution}

\begin{solution}{exo:m1:european-equity-market-structure:5}
$0.07 \times (4\,400 + 400) - 300 = \text{\euro}36$ million, that is 9\% of
next month's volume: a share can run above 7\% in a month as long as the year
stays below.
\end{solution}

\begin{solution}{exo:m1:european-equity-market-structure:6}
Before: two ticks of 0.005, \euro 0.01, that is \qty{8.1}{\bp}. After: the
minimum spread is one tick, again \euro 0.01, \qty{8.1}{\bp}: the quoted
spread need not change, but it is now \emph{one} tick. Prices that used to
sit on two levels each side collapse onto one; the queue at the best price
becomes much longer, time priority becomes valuable, and a \omterm{def:m1:what-a-trading-firm-does:market-maker}{market maker}
earns a full tick per round trip but waits longer for each fill.
\end{solution}

\begin{solution}{exo:m1:european-equity-market-structure:7}
Median month of first breach: 28, and all fifty paths breach. With zero
drift the dark share stays near 4.5\%, its starting level, and no path
breaches: the cap binds on trends, not on noise, because a twelve-month sum
averages the monthly noise away.
\end{solution}

\begin{solution}{exo:m1:european-equity-market-structure:8}
It shows that fills were slightly better than the best quote \emph{on the
three venues the algorithm already uses}: a benchmark built from its own
opportunity set, which it can hardly fail to match. It does not show that a
better price was unavailable on the four other venues or from the \omterm{def:m1:european-equity-market-structure:si}{systematic
internalisers}, which is the best-execution question. Repair: benchmark
against a best bid and offer built from all lit venues with a material share
(or, now, the consolidated tape), state the venue list in the report, and
add the midpoint as a second benchmark, since much of the volume trades
there.
\end{solution}

\begin{solution}{pb:m1:european-equity-market-structure:1}
\textbf{1.} \euro 10\,500 million a month; \euro 787.5 million dark.
\textbf{2.} $H = 0.36 + 4\times0.01 = 0.40$: 2.5 effective operators.
\textbf{3.} $H = 0.1444 + 0.04 + 0.0729 + 0.0056 + 0.0016 + 0.0012 = 0.266$:
3.76 effective mechanisms.
\textbf{4.} Usage exceeds 7\%: the \omterm{def:m1:european-equity-market-structure:waivers}{reference price waiver} is suspended in
this share for three months on every venue.
\textbf{5.} Primary book $38 + 0.2\times7.5\times38/65 = 38.88$; closing
auction 20; MTF lit $27.62$; dark 0; periodic $4 + 3.375 = 7.38$; blocks
$3.5 + 1.875 = 5.38$; total 99.25 (0.75 not traded).
\textbf{6.} 39.2\%, 20.2\%, 27.8\%, 0, 7.4\%, 5.4\%.
\textbf{7.} $+3.4$ points.
\textbf{8.} 3.57: fewer effective mechanisms, since one of six has
disappeared.
\textbf{9.} Dark flow \euro 4.8 million a day; 30\% of it (\euro 1.44 million)
no longer gets the midpoint and pays 2 basis points: \euro 288 a day,
\euro 6\,048 a month in this one share.
\textbf{10.} $9 \times 7.5\%/12 = 5.6\%$.
\textbf{11.} The three empty months stay in the window until they are twelve
months old. Usage is 5.6\% for nine months, then 6.25\%, 6.875\%, and 7.5\%
in the twelfth month after restoration: a cycle of twelve months on, three
months off.
\textbf{12.} 7\%.
\textbf{13.} Below 7\% nobody is ever suspended and the total dark volume
over a cycle is larger ($12\times7 = 84$ against $12\times7.5 = 90$ followed
by three months of nothing, that is 90 over fifteen months, or 72 per twelve).
But each operator gains by taking more share while the waiver lasts, and
would let the others bear the restraint: a commons problem.
\textbf{14.} Participants in the lit market: the legislator's view is that
too much dark trading degrades the public price everyone uses.
\textbf{15.} The institutional investor who loses midpoint executions, and
ultimately the savers behind it; part of the flow also stops trading.
\textbf{16.} Lit: an indicative price and quantity are published before the
uncrossing, and any participant may join. Dark: the publication lasts
milliseconds, reveals no individual order, and the price is usually pinned
inside the reference market's quote, so the information content is close to
that of a midpoint cross.
\textbf{17.} Because displaying a very large order would move the price
against it; the waiver protects the order, and blocks are too lumpy and too
rare to substitute for the lit market.
\textbf{18.} Easier: anything measured against a best bid and offer, such as
question 9's spread saving and any effective-spread statistic. Unchanged:
market shares by venue and mechanism and cap usage, which come from trade
reports, not quotes.
\textbf{19.} \textbf{7.4\%.}
\textbf{20.} Order flow is conserved, and migrates to the closest legal
substitute for whatever a rule has just made expensive.
\end{solution}

\begin{solution}{iq:m1:european-equity-market-structure:1}
No \omterm{def:m1:us-equity-market-structure:opr}{order protection rule} in Europe: \omterm{def:m1:european-equity-market-structure:bestex}{best execution} is a broker-level policy,
not a venue-level routing obligation. No consolidated tape until 2026, so no
official best bid and offer. Many listing exchanges and currencies instead of
one national market, with the primary exchange dominant in its own stocks.
Tick sizes by table (price and liquidity), not a uniform cent. Dark trading
by waiver, with a volume cap; large closing auctions and \omterm{def:m1:european-equity-market-structure:periodic}{periodic auctions}.
No \omterm{def:m1:retail-flow-and-wholesaling:pfof}{payment for order flow} in most jurisdictions; bilateral trading through
\omterm{def:m1:european-equity-market-structure:si}{systematic internalisers} instead of \omterm{def:m1:retail-flow-and-wholesaling:wholesaler}{wholesalers}.
\iqlookfor{the routing/best-execution contrast first; the rest is detail.}
\end{solution}

\begin{solution}{iq:m1:european-equity-market-structure:2}
An investment firm dealing on own account against client orders outside a
venue, on an organised, frequent and substantial basis. Every trade has the
firm as counterparty and the firm chooses its clients; an MTF is a neutral
multilateral system where participants' orders interact under
non-discretionary rules and the operator takes no position.
\iqlookfor{bilateral and principal versus multilateral and neutral.}
\end{solution}

\begin{solution}{iq:m1:european-equity-market-structure:3}
Index funds and benchmarked managers are measured against the close and want
zero \omterm{def:m1:the-buy-side:benchmark}{tracking error} against it; derivatives, funds and indices settle on it;
and the auction offers large size with no spread. As volume concentrates
there, the continuous book thins, intraday impact rises, and still more
participants wait for the close: a self-reinforcing migration that also
hands the listing exchange pricing power over its auction.
\iqlookfor{the feedback loop, not only ``because of passive''.}
\end{solution}

\begin{solution}{iq:m1:european-equity-market-structure:4}
Which venues to include (listing exchange only, all lit venues, those above a
market-share threshold, \omterm{def:m1:european-equity-market-structure:si}{systematic internalisers}' quotes or not); how to
treat periodic-auction indications; currency conversion for multi-listed
shares; clock synchronisation across venues in different data centres; how to
handle a primary-exchange halt or auction while other venues keep trading;
whether to include sizes and what minimum size. None of this is prescribed,
so the choices must be documented and kept fixed across the analysis.
\iqlookfor{the benchmark as a modelling choice that must be declared.}
\end{solution}

\begin{solution}{iq:m1:european-equity-market-structure:5}
Midpoint dark liquidity disappears for three months. Shift the passive,
spread-saving part of the order to \omterm{def:m1:european-equity-market-structure:periodic}{periodic auctions} and to conditional block
venues under the \omterm{def:m1:european-equity-market-structure:waivers}{large-in-scale waiver}; expect lower fill rates at the
midpoint and therefore either a longer schedule or more lit, spread-paying
execution; use the close more. Before the suspension date, if the forecast is
reliable, accelerate the dark-eligible part of the order.
\iqlookfor{substitution toward \omterm{def:m1:european-equity-market-structure:periodic}{periodic auctions} and blocks, and the cost in
fill probability.}
\end{solution}

\begin{solution}{iq:m1:european-equity-market-structure:6}
If the old spread was two or more old ticks, it can now be one new tick: the
quoted spread stays the same or widens to the new minimum; depth at the best
price rises sharply because orders pile onto fewer levels; queue times
lengthen and cancellations fall; price-improvement inside the spread becomes
impossible. For a \omterm{def:m1:what-a-trading-firm-does:market-maker}{market maker}, revenue per fill rises and fills per unit of
time fall, and the advantage shifts from pricing skill to queue position and
speed; trading migrates toward midpoint and periodic venues that escape the
tick. Test with a difference-in-differences on the annual recalculation:
shares that changed band against similar shares that just did not, comparing
spread, depth, \omterm{def:m1:retail-flow-and-wholesaling:effective}{realised spread} and venue shares before and after.
\iqlookfor{queue priority as the new scarce resource, and a clean
identification strategy.}
\end{solution}
